\begin{pastexam}[label={px:ncu-114-dsa-13}]{中央資工所 114 學年度 第 13 題}
（複選；每題 5 分，全部答對才給分，答錯倒扣 1 分）

Which of the following statements about Minimum Spanning Tree are incorrect?
\begin{enumerate}[label=(\alph*)]
\item The spanning tree of a graph exists uniquely which connects all vertices with the minimum total
edge weight.
\item Both Prim's and Kruskal's algorithms are greedy algorithms.
\item After optimization, Prim's and Kruskal's algorithms have the same (worst case) time complexity.
\item Disjoint-set data structure can be used to improve the time complexity of the Prim's algorithm.
\item Fibonacci heap can be used to improve the time complexity of the Kruskal's algorithm.
\end{enumerate}
\end{pastexam}

\begin{pastsoln}
答案：(a)(c)(d)(e)。 題目問的是\textbf{錯誤}的敘述。 以下 $n = |V|$、$m = |E|$。

\begin{itemize}
\item (a) 錯誤。 最小生成樹不一定唯一（資結書第~24.1.7~節）：三角形三條邊權重都是 $1$，任兩條邊都是
最小生成樹，有三棵。 邊權重全部不同時才唯一。 另外圖不連通時根本沒有生成樹。
\item (b) 正確。 兩者每一步都挑目前最便宜、又不違反條件的邊（Kruskal 挑全域最小、不成環的邊；Prim 挑
跨過目前樹的最小邊），正確性都來自切割性質（資結書第~24.1~節）。
\item (c) 錯誤。 Kruskal 要排序邊，$O(m \log n)$；Prim 用費氏堆積可以做到 $O(m + n \log n)$
（資結書第~24.3~節）。 稠密的圖（$m = \Theta(n^2)$）上，前者是 $\Theta(n^2 \log n)$，後者是
$\Theta(n^2)$，不一樣。 只用二元堆積時兩者都是 $O(m \log n)$，但題目說的是「最佳化之後」。
\item (d) 錯誤。 互斥集合（disjoint set）是 \textbf{Kruskal} 用來判斷兩端點是否已經連通、避免成環的
結構（資結書第~24.2~節）。 Prim 只長一棵樹，判斷「在不在樹裡」用一個布林陣列就好，用不到互斥集合。
\item (e) 錯誤。 費氏堆積加快的是 \textbf{Prim}（$n$ 次 \textsc{Extract-Min}、$m$ 次
\textsc{Decrease-Key}，後者平攤 $O(1)$）。 Kruskal 的瓶頸是排序，用不到優先佇列的 \textsc{Decrease-Key}。
\end{itemize}
(d)(e) 是把兩個演算法的資料結構對調了。
\end{pastsoln}
